\documentclass{article}

% NeurIPS 2025 style (vendored in neurips/neurips_2025.sty), preprint mode:
% these are self-study reports, not conference submissions.
\PassOptionsToPackage{numbers,sort&compress}{natbib}
\usepackage[preprint]{neurips/neurips_2025}

\usepackage[T1]{fontenc}
\usepackage{amsmath,amssymb}
\usepackage{booktabs}
\usepackage{graphicx}
\usepackage[colorlinks=true,linkcolor=blue,urlcolor=blue,citecolor=blue]{hyperref}

\title{Attention logit scaling: a controlled empirical mechanism check}
\author{self-study-with-ai}
\date{\today}

\begin{document}
\maketitle

\begin{abstract}
The Transformer scales query-key dot products by $1/\sqrt{d_k}$, and
Vaswani et al.~\citep{vaswani2017attention} motivate the factor by the
suspicion that large dot products push the softmax into regions of small
gradients, but the paper provides no ablation isolating the scale. We check
that stated mechanism
under controlled synthetic sampling. Drawing independent zero-mean Gaussian
query and key coordinates, we measure logit variance, softmax concentration,
and softmax Jacobian norm with and without the scale across
$d_k \in \{8, 64, 512\}$. Scaling holds the logit standard deviation near one
while the unscaled standard deviation grows as $\sqrt{d_k}$, keeps softmax
entropy and max-probability roughly constant while the unscaled softmax
saturates, and prevents the decay of the Jacobian norm toward zero. Relaxing
the unit-variance assumption shifts the normalizing scale to
$\sqrt{d_k}\,\sigma_q\sigma_k$. Boundary: the results establish the mechanism
under idealized assumptions and do not establish the behavior of trained
Transformers.
\end{abstract}

\section{Introduction}

Scaled dot-product attention divides query-key dot products by $\sqrt{d_k}$
before the softmax~\citep{vaswani2017attention}. The stated motivation is that
for large $d_k$ the dot products grow large in magnitude and push the softmax
into regions of extremely small gradients, and that the division counteracts
this effect~\citep{vaswani2017attention}. The source offers this as an
analytical suspicion rather than a measured result, and reports no controlled
comparison of scaled versus unscaled attention, so the mechanism is stated
rather than measured in the primary source (dossier claim CLM-001).

This study supplies that controlled check on synthetic data. We ask whether,
under independent zero-mean Gaussian query and key coordinates, dividing by
$\sqrt{d_k}$ normalizes the logit variance and reduces softmax concentration
and gradient attenuation, as the standard derivation predicts. Contributions,
each tied to a ledgered claim and the run artifact:

\begin{itemize}
  \item We confirm variance normalization: the unscaled logit standard
    deviation grows as $\sqrt{d_k}$ and the scaled standard deviation is held
    near one (CLM-002).
  \item We quantify the reduction in softmax concentration: unscaled max-probability
    rises and entropy falls with $d_k$, while the scaled softmax stays roughly
    constant (CLM-003).
  \item We quantify the prevention of gradient attenuation: the unscaled
    softmax Jacobian norm decays toward zero with $d_k$ while the scaled norm
    stays roughly constant (CLM-004).
  \item We characterize the normalizing scale under relaxed variance:
    $\sqrt{d_k}\,\sigma_q\sigma_k$ (CLM-005).
\end{itemize}

Scope: the study tests the mechanism on idealized synthetic vectors. It does
not train or load a neural network, and it makes no claim about trained
Transformers.

\section{Background}

Scaled dot-product attention maps a query and a set of key-value pairs to an
output as a weighted sum of values, with weights from a softmax over scaled
query-key dot products~\citep{vaswani2017attention}:
\begin{equation}
\mathrm{Attention}(Q, K, V) = \mathrm{softmax}\!\left(\frac{QK^{\top}}{\sqrt{d_k}}\right) V .
\end{equation}

Let each query and key coordinate be independent with zero mean and variances
$\sigma_q^2$ and $\sigma_k^2$. The dot product $q \cdot k = \sum_{i=1}^{d_k}
q_i k_i$ is a sum of $d_k$ independent terms, each with variance
$\sigma_q^2 \sigma_k^2$, so
\begin{equation}
\mathrm{Var}(q \cdot k) = d_k\, \sigma_q^2 \sigma_k^2,
\qquad
\mathrm{std}(q \cdot k) = \sqrt{d_k}\, \sigma_q \sigma_k .
\end{equation}
Dividing by $\sqrt{d_k}$ yields a standard deviation of $\sigma_q\sigma_k$,
independent of $d_k$; under unit variances this is one. This is the
normalization the scaling factor is meant to provide. Large-magnitude logits
concentrate the softmax: as logit spread grows, the max-probability approaches
one, the entropy approaches zero, and the softmax Jacobian, whose entries are
$p_i(\delta_{ij} - p_j)$, shrinks toward zero. The primary source states this
motivation but does not measure it~\citep{vaswani2017attention}.

\section{Methods and Sources}

\paragraph{Sources.} We reuse the single primary source already gathered in
the companion concept study, Vaswani et al.~\citep{vaswani2017attention}, for
framing only: it states the scaling motivation this study checks. The declared
source budget is two; no new sources were retrieved. The primary evidence is
the synthetic run, not the source.

\paragraph{Synthetic protocol.} We sample query and key vectors with
independent zero-mean Gaussian coordinates using
\texttt{numpy.random.default\_rng} at seed $0$, on CPU only. Dot-product
statistics use $20000$ query-key pairs per $d_k$. Softmax statistics use
$2000$ attention rows, each a query attending to $64$ keys. The grid is
$d_k \in \{8, 64, 512\}$, and the relaxation arm varies the coordinate
standard deviation over $\sigma \in \{0.5, 1.0, 2.0\}$. The run is recorded as
manifest RUN-001-full with its config, script, and output under
\texttt{.research/runs/}; the result artifact is
\texttt{experiments/results/full/summary.json}. All reported results are for
the single seed $0$ with no error bars; sampling-noise deviations from the
closed form are at most about $1.2\%$.

\paragraph{Metrics.} For each setting we measure the empirical standard
deviation of the logits; the softmax max-probability and entropy in nats,
averaged over rows; and the Frobenius norm of the softmax Jacobian, computed in
closed form as the square root of
$\sum_i p_i^2(1-p_i)^2 + (\sum_i p_i^2)^2 - \sum_i p_i^4$ and
averaged over rows. The Jacobian norm is a proxy for gradient magnitude, not a
training gradient.

\section{Findings}

\subsection{Scaling normalizes logit variance}

Table~\ref{tab:logit-std} reports the logit standard deviation. The unscaled
standard deviation grows as $\sqrt{d_k}$ and matches the closed form to within
about $1.2$ percent; dividing by $\sqrt{d_k}$ holds it near one, independent of
$d_k$ (CLM-002).

\begin{table}[h]
\centering
\caption{Logit standard deviation versus $d_k$, $20000$ pairs, unit-variance
coordinates. Measured values from RUN-001-full.}
\label{tab:logit-std}
\begin{tabular}{r r r r}
\toprule
$d_k$ & unscaled std & theory $\sqrt{d_k}$ & scaled std \\
\midrule
8 & 2.862 & 2.828 & 1.012 \\
64 & 8.023 & 8.000 & 1.003 \\
512 & 22.570 & 22.627 & 0.997 \\
\bottomrule
\end{tabular}
\end{table}

\subsection{Scaling reduces softmax concentration}

Table~\ref{tab:concentration} reports softmax concentration over $64$ keys. The
unscaled max-probability rises with $d_k$ and the entropy falls, approaching
the saturated limits of one and zero; the scaled softmax stays roughly
constant, with max-probability near $0.107$ and entropy near $3.69$ nats
(CLM-003).

\begin{table}[h]
\centering
\caption{Softmax concentration over $64$ keys, $2000$ rows, unit-variance
coordinates. Measured values from RUN-001-full.}
\label{tab:concentration}
\begin{tabular}{r r r r r}
\toprule
$d_k$ & unscaled max-prob & unscaled entropy & scaled max-prob & scaled entropy \\
\midrule
8 & 0.428 & 2.053 & 0.106 & 3.689 \\
64 & 0.785 & 0.597 & 0.107 & 3.687 \\
512 & 0.923 & 0.193 & 0.107 & 3.684 \\
\bottomrule
\end{tabular}
\end{table}

\subsection{Scaling prevents gradient attenuation}

Table~\ref{tab:gradient} reports the mean Frobenius norm of the softmax
Jacobian. The unscaled norm decays with $d_k$ toward the saturated limit of
zero, while the scaled norm stays roughly constant near $0.18$ (CLM-004). This
is the gradient-attenuation effect the primary source states
qualitatively~\citep{vaswani2017attention}, here measured under controlled
sampling.

\begin{table}[h]
\centering
\caption{Mean Frobenius norm of the softmax Jacobian over $64$ keys, $2000$
rows, unit-variance coordinates. Measured values from RUN-001-full.}
\label{tab:gradient}
\begin{tabular}{r r r}
\toprule
$d_k$ & unscaled Jacobian norm & scaled Jacobian norm \\
\midrule
8 & 0.2946 & 0.1820 \\
64 & 0.2185 & 0.1833 \\
512 & 0.0972 & 0.1837 \\
\bottomrule
\end{tabular}
\end{table}

\subsection{Relaxed variance shifts the normalizing scale}

With coordinate standard deviation $\sigma$, the unscaled logit standard
deviation is $\sqrt{d_k}\,\sigma^2$, so dividing by $\sqrt{d_k}$ alone leaves
$\sigma^2$ rather than one. Table~\ref{tab:relaxation} reports the $d_k = 64$
slice. The scale that yields unit-variance logits is
$\sqrt{d_k}\,\sigma_q\sigma_k$ (CLM-005). The relaxation arm sets
$\sigma_q = \sigma_k = \sigma$; the asymmetric form follows from the
derivation in the Background and reduces to the symmetric instance that is
measured.

\begin{table}[h]
\centering
\caption{Logit standard deviation under relaxed coordinate variance, $d_k=64$,
$20000$ pairs. Measured values from RUN-001-full.}
\label{tab:relaxation}
\begin{tabular}{r r r r r}
\toprule
$\sigma$ & unscaled std & theory $\sqrt{d_k}\,\sigma^2$ & scaled std & theory $\sigma^2$ \\
\midrule
0.5 & 2.001 & 2.000 & 0.250 & 0.250 \\
1.0 & 7.998 & 8.000 & 1.000 & 1.000 \\
2.0 & 32.217 & 32.000 & 4.027 & 4.000 \\
\bottomrule
\end{tabular}
\end{table}

\section{Related Work}

The nearest material is the primary source itself, which states the scaling
motivation as an analytical argument and provides no ablation isolating the
scale factor~\citep{vaswani2017attention}. This study differs by supplying a
controlled synthetic measurement of exactly that stated mechanism, separating
what the source asserts from what is measured under explicit assumptions. A
survey of the broader literature on attention-scaling variants was declared out
of scope by the source budget, so no further sources are compared here.

\section{Limitations}

The evidence does not establish the behavior of trained Transformers. All
results are on independent zero-mean Gaussian coordinates, an idealization that
learned query and key distributions may violate through non-unit variance,
non-zero mean, and coordinate dependence. The relaxation arm covers a bounded
$\sigma$ grid and varies variance only. The Jacobian norm is a proxy for
gradient magnitude, not a measured training gradient. The study checks the
mechanism the source states; it is not a trained-model ablation of scaled
versus unscaled attention, which the source itself omits.

\section{Conclusion}

Under controlled synthetic sampling, the $1/\sqrt{d_k}$ factor does what the
primary source states it is meant to do: it normalizes the logit variance
independently of $d_k$, keeps the softmax from saturating, and prevents the
softmax Jacobian from decaying toward zero. When the unit-variance assumption
is relaxed, the normalizing scale generalizes to $\sqrt{d_k}\,\sigma_q\sigma_k$.
These findings hold within the stated synthetic scope and do not extend to
trained Transformers, which remain an empirical question this study does not
address.

\bibliographystyle{plainnat}
\bibliography{refs}

\end{document}
